\begin{document}

\section{Понятие тензора в механике сплошных сред}

\subsection{Скалярные, векторные и тензорные величины}

В механике сплошных сред физические величины
классифицируются по их поведению при изменении системы координат.

\textit{Скаляр} — величина, не зависящая от ориентации координатной системы
(например, температура или плотность).

\textit{Вектор} — величина, имеющая направление и компоненты,
которые преобразуются линейно при повороте координат
(например, перемещение или скорость).

Однако для описания внутренних сил и деформаций
скалярных и векторных величин недостаточно.
Для этого вводится более общий объект — тензор.

\subsection{Определение тензора}

\textit{Тензором} называется математический объект,
компоненты которого при замене системы координат
преобразуются по тензорному закону.

Для тензора второго ранга:
\begin{equation}
T'_{ij} = a_{ik} a_{jl} T_{kl},
\end{equation}
где $a_{ik}$ — направляющие косинусы,
описывающие поворот координатной системы.

Главное свойство тензора —
\textbf{физическая величина остаётся неизменной},
меняются только её компоненты.

\subsection{Физический смысл тензора}

Тензор возникает тогда,
    когда физическая величина зависит
\textit{одновременно от направления площадки
и от направления действия силы или деформации}.

Классический пример —
тензор напряжений:
на одной и той же площадке
могут действовать силы в различных направлениях,
и все эти эффекты должны быть описаны в рамках единого объекта.

Таким образом, тензор:
\begin{itemize}
\item содержит полную информацию о физическом состоянии;
\item не зависит от выбора координат;
\item позволяет формулировать законы в инвариантной форме.
\end{itemize}

\section{Тензор перемещений и градиент перемещений}

\subsection{Вектор перемещений}

Движение сплошной среды описывается
вектором перемещений
\begin{equation}
\mathbf{u}(\mathbf{x},t) = (u_1, u_2, u_3).
\end{equation}

Сам по себе вектор перемещений
не является тензором второго ранга,
однако его пространственные производные
порождают тензорные величины.

\subsection{Градиент перемещений}

Градиент перемещений определяется как
\begin{equation}
\frac{\partial u_i}{\partial x_j}.
\end{equation}

Этот объект является тензором второго ранга
и может быть разложен на симметричную и антисимметричную части:
\begin{equation}
\frac{\partial u_i}{\partial x_j}
=
\varepsilon_{ij}
+
\omega_{ij},
\end{equation}
где
\begin{equation}
\varepsilon_{ij}
=
\frac{1}{2}
\left(
\frac{\partial u_i}{\partial x_j}
+
\frac{\partial u_j}{\partial x_i}
\right)
\end{equation}
— тензор деформаций,
\begin{equation}
\omega_{ij}
=
\frac{1}{2}
\left(
\frac{\partial u_i}{\partial x_j}
-
\frac{\partial u_j}{\partial x_i}
\right)
\end{equation}
— тензор вращений.

Антисимметричная часть не влияет на напряжения
в линейной теории упругости.

\section{Тензор деформаций и его инварианты}

Тензор малых деформаций $\varepsilon_{ij}$
характеризует локальное изменение формы и объёма
элементарного объёма среды.

Первый инвариант:
\begin{equation}
I_1^\varepsilon = \varepsilon_{kk} = \nabla \cdot \mathbf{u}
\end{equation}
описывает объёмную деформацию.

Девиатор деформаций:
\begin{equation}
\varepsilon'_{ij}
=
\varepsilon_{ij}
-
\frac{1}{3} \varepsilon_{kk} \delta_{ij}
\end{equation}
характеризует изменение формы без изменения объёма.

\section{Тензор напряжений}

Тензор напряжений $\sigma_{ij}$
определяет внутренние силы в среде
и связывает силу, действующую на площадку,
с ориентацией этой площадки:
\begin{equation}
t_i = \sigma_{ij} n_j.
\end{equation}

Из условия равновесия моментов следует симметрия:
\begin{equation}
\sigma_{ij} = \sigma_{ji}.
\end{equation}

\subsection{Инварианты тензора напряжений}

Первый инвариант:
\begin{equation}
I_1^\sigma = \sigma_{kk}
\end{equation}
определяет гидростатическую часть напряжений.

Девиатор напряжений:
\begin{equation}
s_{ij}
=
\sigma_{ij}
-
\frac{1}{3} \sigma_{kk} \delta_{ij}.
\end{equation}

Второй инвариант девиатора:
\begin{equation}
J_2 = \frac{1}{2} s_{ij} s_{ij}
\end{equation}
характеризует интенсивность касательных напряжений.

\section{Роль тензоров и инвариантов в волновых задачах}

В продольных волнах
определяющим является первый инвариант деформаций
\begin{equation}
\varepsilon_{kk},
\end{equation}
так как именно он отвечает за объёмные изменения среды.

Связь между тензорами деформаций и напряжений
приводит к распространению упругих и термоупругих волн,
скорость которых определяется
объёмными упругими характеристиками материала.

\section{Заключение}

Тензорный аппарат является фундаментальной основой
механики сплошных сред.
Использование тензоров позволяет
формулировать физические законы
в координатно-независимой форме
и корректно описывать волновые процессы
в упругих и термоупругих средах.

\begin{thebibliography}{99}

\bibitem{Landau}
Ландау Л.~Д., Лифшиц Е.~М.
\textit{Теория упругости}.
— М.: Наука, 1987.

\bibitem{Novozhilov}
Новожилов В.~В.
\textit{Основы нелинейной теории упругости}.
— Л.: Судпромгиз, 1948.

\bibitem{Achenbach}
Achenbach J.~D.
\textit{Wave Propagation in Elastic Solids}.
— Amsterdam: North-Holland, 1973.

\bibitem{Nowacki}
Nowacki W.
\textit{Thermoelasticity}.
— Oxford: Pergamon Press, 1986.

\end{thebibliography}

\end{document}
